\worksheettitle
  {Домашняя работа: ответы}
  {\modulemark{M05-H} Версия учителя}

\begin{teacherbox}[Назначение]
  Работа проверяет точное чтение условия, систематический перебор, обязательные
  цифры, ноль в начале записи и обнаружение ошибок в списке.
\end{teacherbox}

\section{Ответы}

\begin{teacherbox}[1. Составьте списки]
  \begin{enumerate}
    \item \(22,\ 26,\ 62,\ 66\);
    \item \(26,\ 28,\ 62,\ 68,\ 82,\ 86\);
    \item \(309,\ 390,\ 903,\ 930\).
  \end{enumerate}
\end{teacherbox}

\begin{teacherbox}[2. Обязательные цифры]
  \(447,\ 474,\ 477,\ 744,\ 747,\ 774\). Не вошли \(444\) и \(777\), потому
  что в каждом из них используется только одна заданная цифра.
\end{teacherbox}

\begin{teacherbox}[3. Доказательство]
  Принимается объяснение, в котором названы все допустимые первые цифры, для
  каждой проверены разрешённые продолжения, а полученное количество совпадает
  с ожидаемым.
\end{teacherbox}

\newpage
\begin{teacherbox}[4. Исправьте список]
  Повторено \(81\), пропущено \(84\). Исправленный список:
  \(14,\ 18,\ 41,\ 48,\ 81,\ 84\).
\end{teacherbox}

\begin{teacherbox}[5. Ноль в начале]
  Записи \(05\) и \(08\) обозначают числа \(5\) и \(8\), поэтому не являются
  двузначными. Полный список: \(50,\ 58,\ 80,\ 85\).
\end{teacherbox}

\begin{teacherbox}[6. Составьте условие]
  Возможная формулировка: «Из цифр \(2\) и \(5\) составьте все двузначные
  числа. Повторения разрешены, использовать обе цифры необязательно».

  Если обе цифры обязательны, остаются \(25,\ 52\).
\end{teacherbox}

\begin{teacherbox}[7. Итог]
  \(207,\ 270,\ 702,\ 720\). Допустимы только первые цифры \(2\) и \(7\);
  после выбора первой цифры две оставшиеся располагаются двумя способами.
\end{teacherbox}

\begin{resultbox}[Критерии проверки]
  \begin{itemize}
    \item \textbf{Освоено:} списки полны и упорядочены, условия не изменены,
      объяснение связано со структурой перебора.
    \item \textbf{Точечная доработка:} одна повторяющаяся ошибка; дать два
      аналогичных коротких списка.
    \item \textbf{M05-R:} случайный перебор, самовольный запрет повторений,
      потеря обязательной цифры или ноль в начале записи.
  \end{itemize}
\end{resultbox}
